\telubab{1}{BAB I}{INTRODUCTION}{sec:1}

\subsection{Background (Latar Belakang)}\label{sec:1.1}

Large language models (LLMs) can answer a wide range of questions, but
their knowledge is fixed when training ends, and they sometimes state
false information with confidence. A common remedy is
\textbf{retrieval-augmented generation} (RAG): before answering, the
system retrieves relevant documents and conditions its answer on them.
Many questions cannot be answered with a single retrieval. A
\emph{multi-hop} question such as ``Which company formed the boy group
that debuted with the album \emph{2014 S/S}?'' requires finding one fact
(the group), then using it to find another (the company).

This has motivated \textbf{search agents}: LLMs that decide for
themselves when to search, what to search for, and when to stop and
answer (Yao et al., 2023). Prompting alone does not reliably produce
good search behaviour, so recent work trains search agents with
\textbf{reinforcement learning} (RL). In particular, \textbf{Search-R1}
(Jin et al., 2025) trains a model with an outcome reward, meaning the
model is rewarded only when its final answer is correct. It uses
\textbf{GRPO} (Shao et al., 2024), an RL algorithm that compares several
attempts at the same question instead of training a separate value
model. GRPO with rule-checked rewards has also been central to recent
reasoning models (DeepSeek-AI, 2025).

Rewarding only the final answer creates a problem. A correct answer does
not show that the agent found, used, or could point to the evidence
behind it. In RL terms, the answer reward is a \textbf{proxy} for what
is actually wanted: a correct answer \emph{supported by evidence}.
Optimizing a proxy can produce behaviour that scores well on the proxy
but not on the real goal (Skalse et al., 2022).

HotpotQA (Yang et al., 2018) offers a way to study this cheaply. Its
questions are annotated with the exact sentences that support the
answer. If an agent cites sentences, the citations can be checked by a
simple rule against these gold labels, with no judge.

\subsection{Problem Statement (Rumusan Masalah)}\label{sec:1.2}

This thesis addresses four related gaps:

\begin{itemize}
\tightlist
\item
  \textbf{G1, correct is not grounded.} An agent trained with
  answer-only reward can learn to answer without showing evidence, as
  long as the final answer is right. Recent work reports such shortcut
  behaviour in outcome-trained search agents (Zhang et al., 2026).
\item
  \textbf{G2, judge-free evidence rewards are rarely tested against an
  outcome-only control.} Several evidence rewards use an LLM as a judge
  (Zhang et al., 2026; Xu et al., 2026), which adds cost and possible
  bias. Rule-based alternatives exist (Chapter~\ref{sec:2}), but they
  are rarely compared over several seeds with an outcome-only reward in
  an otherwise identical setup.
\item
  \textbf{G3, a compute barrier.} RL for search agents is reported at
  multi-GPU scale: Search-R1 trains on 8 H100 GPUs and AutoRefine on 8
  A100 GPUs (Chapter~\ref{sec:2}). It is unclear whether the same
  effects can be studied on one GPU.
\item
  \textbf{G4, weak evidence about what matters.} RL results vary
  considerably between random seeds (Henderson et al., 2018). Search-R1
  and CaRR train each configuration once (Chapter~\ref{sec:2}), and
  single runs cannot separate the effect of a method from seed noise
  (Agarwal et al., 2021).
\end{itemize}

These gaps lead to three research questions:

\begin{itemize}
\tightlist
\item
  \textbf{RQ1.} Does GRPO training improve a 7B search agent on HotpotQA
  over the un-tuned model?
\item
  \textbf{RQ2.} Does a verifiable evidence-grounded reward (VEGR)
  prevent the shortcut behaviour of answer-only reward, without losing
  accuracy?
\item
  \textbf{RQ3.} How do reward variants (gated penalties, a larger
  coverage weight) and optimizer changes (advantages without std
  normalisation, dynamic sampling) affect grounding and search
  behaviour?
\end{itemize}

RQ1 and RQ2 are the primary questions. RQ3 is exploratory: its variants
were partly designed after the first results were seen
(Section~\ref{sec:3.8}).

\subsection{Objectives (Tujuan)}\label{sec:1.3}

Each objective answers one research question:

\begin{enumerate}
\def\labelenumi{\arabic{enumi}.}
\tightlist
\item
  \textbf{For RQ1:} train a search agent with GRPO-style RL and LoRA on
  a single GPU, and measure its accuracy and protocol compliance against
  the un-tuned model on a frozen holdout.
\item
  \textbf{For RQ2:} design a judge-free reward, VEGR, that rewards
  correct answers, correct citations and useful searches, and compare it
  with an answer-only reward over three seeds, on both accuracy and
  grounding, with seed-aware bootstrap confidence intervals.
\item
  \textbf{For RQ3:} test whether reward variants and optimizer changes
  alter grounding, coverage and training stability.
\end{enumerate}

\subsection{Scope and Limitations (Batasan Masalah)}\label{sec:1.4}

This thesis is deliberately narrow:

\begin{itemize}
\tightlist
\item
  \textbf{Data and retrieval.} It uses HotpotQA in the distractor
  setting. The agent searches \textbf{a closed pool of each question's
  10 paragraphs} with BM25, and this pool always contains both gold
  paragraphs. The un-tuned model already retrieves 83\% of the gold
  documents (retrieval recall 0.83), so this is much easier than
  open-web or open-corpus search, and ``search agent'' in this thesis
  means an agent searching this closed pool.
\item
  \textbf{Model.} The model is Qwen2.5-7B-Instruct adapted with LoRA.
\item
  \textbf{Scale.} Each training run uses 256 questions for one epoch,
  and evaluation uses a frozen 300-question holdout. The project was
  budgeted at about \$45; the actual compute spend was about \$49
  (Table~\ref{tbl:3.2}).
\item
  \textbf{Labels.} VEGR needs sentence-level gold evidence, which few
  datasets provide. The method cannot be applied as-is where such labels
  do not exist.
\item
  \textbf{What is not claimed.} The thesis does not propose a new RL
  algorithm or a new kind of reward, does not claim state-of-the-art
  accuracy, and does not claim that the findings transfer to open-web
  search, other datasets or larger models.
\item
  \textbf{Known loophole.} The environment accepts an answer with no
  citations. This was kept on purpose so that all runs are comparable,
  and its consequences are reported.
\end{itemize}

Chapter~\ref{sec:5} discusses the further limitations.

\subsection{Benefits (Manfaat)}\label{sec:1.5}

\textbf{Contributions.}

\begin{enumerate}
\def\labelenumi{\arabic{enumi}.}
\tightlist
\item
  \textbf{VEGR}, a judge-free, sentence-level reward for training search
  agents. Every term is checked by rule against gold supporting facts.
  VEGR also enforces the action protocol: invented citations, malformed
  actions, and running out of budget all receive a −1 penalty. The
  individual reward terms are not new; the contribution is the
  combination and the controlled evaluation of it.
\item
  \textbf{Evidence of when answer-only reward loses citations, and why.}
  Two of three outcome-only seeds stopped citing entirely
  (Clopper--Pearson 95\% CI for the rate {[}0.09, 0.99{]}). Training
  logs show why: under the answer-only reward, invalid citations cost −1
  and correct ones earn nothing. The full VEGR bundle reached the same
  accuracy (+1.0 EM points, seed-aware 95\% CI {[}−3.6, +5.4{]}) but
  \textbf{24.4 points higher precision-grounded EM} (CI {[}+8.8,
  +37.3{]}).
\item
  \textbf{A controlled multi-seed comparison} with seed-aware bootstrap
  confidence intervals. Exploratory results show that:

  \begin{itemize}
  \tightlist
  \item
    a single-seed improvement of +6 EM points (gated penalties, R2) did
    not replicate across seeds;
  \item
    VEGR makes citations correct but not complete, and a larger coverage
    weight (R3) trades precision for completeness, ranking best on
    HotpotQA's Joint-EM;
  \item
    on the one seed tested, dynamic sampling raised outcome-only
    accuracy (together with twice the training questions) but did not
    prevent the loss of citations.
  \end{itemize}
\item
  \textbf{A documented single-GPU pipeline.} One training run plus
  evaluation costs about \$2.25 on one 48 GB GPU. Splits, configs,
  training logs and per-episode results are in the repository; LoRA
  adapters are kept for 15 of the 16 runs.
\end{enumerate}

\textbf{Benefits.} For researchers, the thesis shows how an evidence
reward can be evaluated cheaply and with honest uncertainty, and it
gives a concrete account of how answer-only rewards can lose citations.
For practitioners building retrieval-augmented systems, it shows that
answer accuracy alone does not reveal whether answers are grounded, and
that grounding has to be rewarded and measured explicitly.

\subsection{Thesis Structure}\label{sec:1.6}

\begin{itemize}
\tightlist
\item
  \textbf{Chapter~\ref{sec:2}} reviews the background: LLMs, retrieval,
  multi-hop QA and citation evaluation, LLM agents, RL for LLMs (PPO,
  GRPO and its variants, LoRA), RL-trained search agents, evidence
  rewards, and reward hacking and RL evaluation. It ends with the
  research gaps.
\item
  \textbf{Chapter~\ref{sec:3}} describes the methodology: data splits,
  the search environment, the VEGR reward and its variants, the policy
  model, the training algorithm, the evaluation protocol, and the
  experiment matrix.
\item
  \textbf{Chapter~\ref{sec:4}} reports the results by research question,
  diagnoses the training dynamics, presents example traces, and
  discusses threats to validity.
\item
  \textbf{Chapter~\ref{sec:5}} concludes, answers the research
  questions, and outlines ethics, limitations and future work.
\end{itemize}
